\documentclass[11pt,a4paper]{article}

\usepackage[T1]{fontenc}
\usepackage[utf8]{inputenc}
\usepackage[margin=25mm]{geometry}
\usepackage{amsmath,amssymb,bm}
\usepackage{booktabs}
\usepackage{array}
\usepackage{graphicx}
\usepackage{enumitem}
\usepackage{tikz}
\usetikzlibrary{circuits.ee.IEC,positioning,calc}
\usepackage[hidelinks]{hyperref}

\setlength{\emergencystretch}{3em}
\newcommand{\ohm}{\,\Omega}
\newcommand{\ud}{\,\mathrm{d}}
\newcommand{\Zmat}{\bm{Z}}
\newcommand{\Ymat}{\bm{Y}}
\newcommand{\Rmat}{\bm{R}}
\newcommand{\Lmat}{\bm{L}}
\newcommand{\Cmat}{\bm{C}}
\newcommand{\Gmat}{\bm{G}}
\newcommand{\Vvec}{\bm{V}}
\newcommand{\Ivec}{\bm{I}}

\title{\textbf{Mathematical Model and Lumped RLC Formulation for the\\
Simscape Implementation of a Shielded Three-Phase Cable}}
\author{Davide Bagnara}
\date{\today}

\begin{document}
\maketitle

\begin{abstract}
\noindent
This document develops the mathematical model required to implement a shielded
three-phase power cable in Simscape (MATLAB) as a lumped RLC network. The
multiconductor transmission-line (MTL) formulation is reduced to per-unit-length
(p.u.l.) series and shunt matrices, then mapped onto single- and multi-section
$\pi$ models. A complete numerical dimensioning is carried out for the reference
case of a copper cable with $90\ \mathrm{mm^2}$ per phase, collectively shielded,
of length $120\ \mathrm{m}$. The section-count criterion for the cascaded $\pi$
model is derived from the propagation velocity and the target bandwidth, and the
shield/common-mode representation is included so that the same model serves both
fundamental-frequency studies and high-frequency phenomena (reflected-wave
overvoltage, common-mode current). References follow IEEE practice.
\end{abstract}

\tableofcontents
\bigskip

%======================================================================
\section{Scope and modeling objectives}
%======================================================================
The objective is a circuit-equivalent cable model suitable for time-domain
simulation in Simscape / Simscape Electrical (Specialized Power Systems). Three
requirements drive the formulation:

\begin{enumerate}[leftmargin=*]
\item \textbf{Fidelity at the fundamental and low harmonics} for steady-state and
power-quality studies (voltage drop, charging current, losses).
\item \textbf{Representation of the metallic shield} as an explicit conductor, so
that phase-to-shield capacitance and shield bonding can be analyzed, which is the
prerequisite for common-mode and EMC studies.
\item \textbf{Controlled bandwidth} through a cascaded $\pi$ structure whose number
of sections is chosen for the highest frequency of interest, allowing the same
topology to capture cable resonances and reflected-wave behavior at PWM-converter
output.
\end{enumerate}

The cable is treated as a uniform, laterally homogeneous line. End effects,
joints and terminations are excluded; they are added as separate lumped elements
when needed.

%======================================================================
\section{Cable construction and reference data}
%======================================================================
The reference cable is a three-core copper construction with a collective metallic
shield, representative of a $0.6/1\ \mathrm{kV}$ EMC drive cable. The assumed
geometry is summarized in Table~\ref{tab:geom}. These values define the worked
example; for a specific product the manufacturer datasheet p.u.l. values should be
substituted directly into Sections~\ref{sec:numbers}--\ref{sec:simscape}.

\begin{table}[h]
\centering
\caption{Reference geometry and material data.}
\label{tab:geom}
\begin{tabular}{@{}lll@{}}
\toprule
Quantity & Symbol & Value \\
\midrule
Phase conductor cross-section & $A_c$ & $90\ \mathrm{mm^2}$ (Cu, stranded) \\
Effective conductor radius & $r_c$ & $6.0\ \mathrm{mm}$ \\
Geometric mean radius & $\mathrm{GMR}$ & $0.77\,r_c \approx 4.6\ \mathrm{mm}$ \\
Insulation (XLPE) rel. permittivity & $\varepsilon_r$ & $2.5$ \\
Insulation loss tangent & $\tan\delta$ & $4\times10^{-3}$ \\
Radius over core insulation & $r_i$ & $8.0\ \mathrm{mm}$ \\
Core lay-up & --- & trefoil, touching \\
Center-to-center spacing & $S$ & $16\ \mathrm{mm}$ \\
Shield inner radius & $R_s$ & $18\ \mathrm{mm}$ \\
Shield cross-section (Cu) & $A_s$ & $16\ \mathrm{mm^2}$ \\
Cable length & $\ell$ & $120\ \mathrm{m}$ \\
Copper resistivity ($20\,^\circ\mathrm{C}$) & $\rho_{20}$ & $1.724\times10^{-8}\ \Omega\,\mathrm{m}$ \\
Temperature coefficient & $\alpha$ & $3.93\times10^{-3}\ \mathrm{K^{-1}}$ \\
\bottomrule
\end{tabular}
\end{table}

%======================================================================
\section{Distributed-parameter formulation}
%======================================================================
With phase conductors $a,b,c$ and shield $s$, the cable is a four-conductor MTL.
In the frequency domain the telegrapher equations are
\begin{equation}
-\frac{\ud \Vvec}{\ud x}=\Zmat\,\Ivec,
\qquad
-\frac{\ud \Ivec}{\ud x}=\Ymat\,\Vvec,
\label{eq:telegraph}
\end{equation}
with the p.u.l. series impedance and shunt admittance matrices
\begin{equation}
\Zmat=\Rmat+j\omega\Lmat,
\qquad
\Ymat=\Gmat+j\omega\Cmat .
\label{eq:ZY}
\end{equation}
$\Vvec$ collects conductor-to-reference voltages and $\Ivec$ the conductor
currents. For a symmetric three-phase cable the $3\times3$ phase blocks have the
balanced structure
\begin{equation}
\Zmat_{abc}=
\begin{bmatrix}
Z_s & Z_m & Z_m\\
Z_m & Z_s & Z_m\\
Z_m & Z_m & Z_s
\end{bmatrix},
\qquad
\Ymat_{abc}=
\begin{bmatrix}
Y_s & Y_m & Y_m\\
Y_m & Y_s & Y_m\\
Y_m & Y_m & Y_s
\end{bmatrix}.
\label{eq:balanced}
\end{equation}

\paragraph{Shield elimination (Kron reduction).}
When the shield is bonded to ground at both ends, its voltage is constrained and
the four-conductor system reduces to an effective $3\times3$ phase model:
\begin{equation}
\Zmat_{abc}^{\,\mathrm{eff}}
=\Zmat_{pp}-\Zmat_{ps}\,Z_{ss}^{-1}\,\Zmat_{sp},
\label{eq:kron}
\end{equation}
where $pp$, $ps$, $sp$, $ss$ are the phase, phase-shield, shield-phase and shield
partitions of the full $4\times4$ matrix. Single-point bonding suppresses the
shield circulating current and is treated by keeping the shield as an explicit
floating conductor instead of applying \eqref{eq:kron}. The lumped model below
retains the shield explicitly so that either bonding scheme can be imposed in
Simscape.

\paragraph{Modal / sequence decomposition.}
The balanced structure \eqref{eq:balanced} is diagonalized by the symmetrical-component
(or, in real arithmetic, Clarke) transform, giving decoupled positive/negative
and zero modes
\begin{equation}
Z_1=Z_s-Z_m,\qquad Z_0=Z_s+2Z_m,
\qquad
Y_1=Y_s-Y_m,\qquad Y_0=Y_s+2Y_m .
\label{eq:seq}
\end{equation}
Equation~\eqref{eq:seq} is the bridge between the physical p.u.l. quantities and
the positive-/zero-sequence parameters used by the Simscape three-phase line
block.

%======================================================================
\section{Per-unit-length parameters}\label{sec:pul}
%======================================================================

\subsection{Series resistance}
The DC resistance of one conductor is
\begin{equation}
R'_{\mathrm{dc}}(\theta)=\frac{\rho_{20}}{A_c}\bigl[1+\alpha(\theta-20)\bigr].
\label{eq:rdc}
\end{equation}
The AC resistance adds skin ($y_s$) and proximity ($y_p$) corrections following
IEC~60287 / IEEE cable-loss practice:
\begin{equation}
R'_{\mathrm{ac}}=R'_{\mathrm{dc}}\,(1+y_s+y_p),
\qquad
y_s=\frac{x_s^4}{192+0.8\,x_s^4},
\qquad
x_s^2=\frac{8\pi f}{R'_{\mathrm{dc}}}\,10^{-7}\,k_s,
\label{eq:rac}
\end{equation}
with $k_s\approx1$ for round stranded conductors. At $50\ \mathrm{Hz}$ the
correction is a few percent for $90\ \mathrm{mm^2}$; at the kHz range relevant to
converter harmonics it grows and is captured by re-evaluating \eqref{eq:rac} at the
dominant harmonic, or by a Foster (R--L ladder) branch when broadband fidelity is
required.

\subsection{Series inductance}
The positive-sequence (loop) inductance per phase is
\begin{equation}
L_1'=\frac{\mu_0}{2\pi}\ln\!\frac{\mathrm{GMD}}{\mathrm{GMR}},
\qquad \mathrm{GMD}=\sqrt[3]{S_{ab}S_{bc}S_{ca}}=S \;\;(\text{trefoil}).
\label{eq:l1}
\end{equation}
The zero-sequence inductance depends on the common-mode return path. With the
shield as the dominant return at radius $R_s$,
\begin{equation}
L_0'\approx\frac{\mu_0}{2\pi}\ln\!\frac{R_s}{\mathrm{GMR}}
\;+\;\Delta L_{\mathrm{earth}},
\label{eq:l0}
\end{equation}
where $\Delta L_{\mathrm{earth}}$ accounts for the earth-return contribution
(Carson/Pollaczek) when the shield is grounded at both ends. $L_0'$ is therefore
bonding-dependent and must be re-evaluated for the actual installation.

\subsection{Shunt capacitance}
Two physical capacitances describe the shielded cable: the core-to-shield
capacitance $C_g'$ and the core-to-core capacitance $C_m'$. For an individually
screened core the former is the coaxial value
\begin{equation}
C_g'=\frac{2\pi\varepsilon_0\varepsilon_r}{\ln(r_i/r_c)} ,
\label{eq:cg}
\end{equation}
while for a collective shield $C_g'$ and $C_m'$ are obtained from the Maxwell
potential-coefficient matrix $\bm{P}=\Cmat^{-1}/(2\pi\varepsilon)$ of the
core/shield geometry. The sequence capacitances follow from \eqref{eq:seq}:
\begin{equation}
C_1'=C_g'+3\,C_m', \qquad C_0'=C_g'.
\label{eq:cseq}
\end{equation}
$C_1'$ is the operating (line) capacitance that sets the charging current; $C_0'$
is the capacitance seen by the common mode.

\subsection{Shunt conductance}
Dielectric losses give
\begin{equation}
G'=\omega\,C'\tan\delta,
\label{eq:g}
\end{equation}
negligible at $50\ \mathrm{Hz}$ but retained as a small parallel conductance to
damp the high-frequency $\pi$-network resonances and avoid undamped numerical
oscillation.

\subsection{Shield parameters}
The shield DC resistance is $R_s'=\rho_{20}/A_s$; at high frequency the braid
contributes a transfer impedance, but for the lumped model the shield is
represented by its series $R_s'$, its self-inductance, and the $C_g'$ coupling to
each core.

%======================================================================
\section{Numerical dimensioning of the reference cable}\label{sec:numbers}
%======================================================================
Applying \eqref{eq:rdc}--\eqref{eq:g} to Table~\ref{tab:geom} gives the p.u.l.
values in Table~\ref{tab:pul}. The conductor AC resistance is reported at the
$90\,^\circ\mathrm{C}$ rated temperature with a $50\ \mathrm{Hz}$ skin/proximity
correction of about $2\%$.

\begin{table}[h]
\centering
\caption{Computed per-unit-length parameters of the reference cable.}
\label{tab:pul}
\begin{tabular}{@{}lcrl@{}}
\toprule
Parameter & Symbol & Value & Unit \\
\midrule
Conductor resistance, DC $20\,^\circ\mathrm{C}$ & $R'_{\mathrm{dc}}$ & $0.192$ & $\Omega/\mathrm{km}$ \\
Conductor resistance, AC $90\,^\circ\mathrm{C}$, $50\,\mathrm{Hz}$ & $R'_{\mathrm{ac}}$ & $0.248$ & $\Omega/\mathrm{km}$ \\
Shield resistance, DC $20\,^\circ\mathrm{C}$ & $R'_{s}$ & $1.08$ & $\Omega/\mathrm{km}$ \\
Positive-sequence inductance & $L_1'$ & $0.25$ & $\mathrm{mH/km}$ \\
Zero-sequence inductance (both-end bonding) & $L_0'$ & $0.27$ & $\mathrm{mH/km}$ \\
Core-to-shield capacitance & $C_g'$ & $0.30$ & $\mu\mathrm{F/km}$ \\
Core-to-core capacitance & $C_m'$ & $0.10$ & $\mu\mathrm{F/km}$ \\
Positive-sequence capacitance & $C_1'=C_g'+3C_m'$ & $0.60$ & $\mu\mathrm{F/km}$ \\
Zero-sequence capacitance & $C_0'=C_g'$ & $0.30$ & $\mu\mathrm{F/km}$ \\
Shunt conductance ($50\,\mathrm{Hz}$) & $G_1'$ & $0.75$ & $\mu\mathrm{S/km}$ \\
\bottomrule
\end{tabular}
\end{table}

\paragraph{Lumped totals for $\ell=120\ \mathrm{m}$.}
Multiplying by $\ell$ gives the total elements of the equivalent network:
\begin{align*}
R &= R'_{\mathrm{ac}}\,\ell = 0.248\times0.120 = 29.8\ \mathrm{m}\Omega \ \text{per phase},\\
L_1 &= L_1'\,\ell = 30.0\ \mu\mathrm{H} \ \text{per phase},\qquad
L_0 = 32.4\ \mu\mathrm{H},\\
C_1 &= C_1'\,\ell = 72\ \mathrm{nF} \ \text{(operating, line)},\qquad
C_g = C_g'\,\ell = 36\ \mathrm{nF} \ \text{per phase to shield},\\
R_s &= R_s'\,\ell = 130\ \mathrm{m}\Omega \ \text{(shield)}.
\end{align*}

\paragraph{Derived line quantities.}
The propagation velocity, characteristic impedance and one-way delay of the
positive mode are
\begin{equation}
v=\frac{1}{\sqrt{L_1'C_1'}}=8.2\times10^{7}\ \mathrm{m/s},
\qquad
Z_0=\sqrt{\frac{L_1'}{C_1'}}\approx20\ \Omega,
\qquad
\tau=\frac{\ell}{v}=1.47\ \mu\mathrm{s}.
\label{eq:derived}
\end{equation}
The first open-end (quarter-wave) resonance of the $120\ \mathrm{m}$ cable is
\begin{equation}
f_{\lambda/4}=\frac{v}{4\ell}\approx170\ \mathrm{kHz},
\label{eq:fres}
\end{equation}
which sets the bandwidth the lumped model must reach to reproduce the first
resonance and the reflected-wave behavior at a converter output.

%======================================================================
\section{Lumped RLC model}\label{sec:lumped}
%======================================================================

\subsection{Single $\pi$ section}
For each phase the elementary nominal-$\pi$ cell places the total series branch
$R+j\omega L$ between two halves of the total shunt branch, $Y/2$ at each end.
Including the shield explicitly, the shunt branch is split into the core-to-shield
capacitance $C_g$ and the core-to-core capacitance $C_m$, and the shield node
carries $R_s,L_s$ to its bonding point. The single-$\pi$ cell is accurate while the
cable is electrically short, i.e. $\ell \ll \lambda/10$. From \eqref{eq:derived}
this holds up to
\begin{equation}
f_{\max}^{(1\pi)}=\frac{v}{10\,\ell}\approx 68\ \mathrm{kHz},
\end{equation}
adequate for fundamental and low-harmonic studies but not for the resonance in
\eqref{eq:fres}.

\subsection{Cascaded $N$-section $\pi$ model}
For broadband fidelity the cable is split into $N$ identical $\pi$ cells, each
representing a length $\ell/N$ with element values $R/N$, $L/N$, $C/N$, $G/N$. The
number of cells is set by the highest frequency of interest $f_{\max}$ using a
``cells-per-wavelength'' rule. Requiring $k$ cells per shortest wavelength
$\lambda_{\min}=v/f_{\max}$,
\begin{equation}
\boxed{\,N \ge k\,\frac{\ell\,f_{\max}}{v}\,},\qquad k\approx 8 .
\label{eq:Ncriterion}
\end{equation}
For the reference cable, Table~\ref{tab:N} lists the resulting section count.

\begin{table}[h]
\centering
\caption{Required number of $\pi$ sections for the $120\ \mathrm{m}$ reference cable
($k=8$).}
\label{tab:N}
\begin{tabular}{@{}lcc@{}}
\toprule
Target bandwidth $f_{\max}$ & $\lambda_{\min}=v/f_{\max}$ & $N$ \\
\midrule
$10\ \mathrm{kHz}$  & $8.2\ \mathrm{km}$  & $1$ \\
$100\ \mathrm{kHz}$ & $820\ \mathrm{m}$   & $2$ \\
$500\ \mathrm{kHz}$ & $164\ \mathrm{m}$   & $6$ \\
$1\ \mathrm{MHz}$   & $82\ \mathrm{m}$    & $12$ \\
$3\ \mathrm{MHz}$   & $27\ \mathrm{m}$    & $35$ \\
\bottomrule
\end{tabular}
\end{table}

\noindent
A practical choice for PWM-drive studies (reflected wave, $dv/dt$ stress) is
$N=20$--$40$, since the relevant edge spectra extend into the low-MHz range. The
trade-off is computational: each added section introduces one more $L$--$C$ state
pair and lowers the simulation step required for stability.

\subsection{Phase-domain versus sequence-domain realization}
Two equivalent realizations of each cell are possible:
\begin{itemize}[leftmargin=*]
\item \textbf{Sequence-domain:} use $R_1,L_1,C_1$ and $R_0,L_0,C_0$ from
Table~\ref{tab:pul}. Valid for the balanced, transposed assumption and directly
accepted by the Simscape three-phase $\pi$-line block.
\item \textbf{Phase-domain:} populate the $3\times3$ (or $4\times4$ with shield)
$\Rmat,\Lmat,\Cmat$ matrices from \eqref{eq:balanced}. This retains explicit
inter-phase and phase-shield coupling and is preferred for unbalanced or
common-mode studies.
\end{itemize}

\subsection{Shield and common-mode representation}
The common-mode loop is the three phases in parallel returning through the shield
(and earth). Its lumped parameters are obtained from the zero-sequence quantities:
the equivalent CM inductance is $L_0/3$ in series with the shield $R_s,L_s$, and
the CM capacitance is $3C_g$. Retaining the shield node with selectable bonding
(both-end, single-end, or floating) lets the same model reproduce shield
circulating current and the high-frequency CM path that governs bearing current
and conducted EMI in drive applications.

%======================================================================
\section{Simscape implementation}\label{sec:simscape}
%======================================================================
The model maps directly onto Simscape Electrical components.

\begin{enumerate}[leftmargin=*]
\item \textbf{Ready-made block.} The \emph{Three-Phase PI Section Line} (Specialized
Power Systems) accepts the positive- and zero-sequence $R_1,R_0$
($\Omega/\mathrm{km}$), $L_1,L_0$ ($\mathrm{H/km}$), $C_1,C_0$
($\mathrm{F/km}$), the length, and the number of sections $N$. Enter the values
of Table~\ref{tab:pul}, $\ell=0.12\ \mathrm{km}$, and $N$ from
Table~\ref{tab:N}. This is the fastest route when an explicit shield node is not
needed.
\item \textbf{Explicit shield / custom cell.} When the shield must be accessible,
build one $\pi$ cell from foundation components --- \emph{Resistor},
\emph{Inductor} (or mutual-inductor / \emph{Three-Phase Mutual Inductance Z1-Z0}
for coupling), and \emph{Capacitor} --- with the per-cell values $R/N$, $L/N$,
$C_g/N$, $C_m/N$, $G/N$, and a dedicated shield node carrying $R_s/N$, $L_s/N$.
Cascade $N$ cells and expose the shield terminal for bonding.
\item \textbf{Numerical settings.} The cascaded $\pi$ network is stiff; use a
variable-step stiff solver (\texttt{ode23t}/\texttt{ode15s}) or a fixed-step
discrete solver with $T_s \lesssim 1/(20 f_{\max})$. The small dielectric
conductance $G$ from \eqref{eq:g} damps the inter-cell resonances and improves
conditioning.
\item \textbf{Parameterization.} Drive the block masks from a MATLAB script that
computes Table~\ref{tab:pul} from the geometry of Table~\ref{tab:geom}, so the same
model regenerates for any length or cross-section.
\end{enumerate}

%======================================================================
\section{Validation}
%======================================================================
The lumped model is validated against the closed-form distributed line over the
intended bandwidth:
\begin{enumerate}[leftmargin=*]
\item \textbf{Input impedance.} Compare the open- and short-circuit driving-point
impedance of the $N$-cell network with the exact
$Z_{\mathrm{in}}=Z_0\coth(\gamma\ell)$ and $Z_0\tanh(\gamma\ell)$, where
$\gamma=\sqrt{\Zmat\Ymat}$. The first resonance must match \eqref{eq:fres} within a
few percent.
\item \textbf{Charging current.} At $50\ \mathrm{Hz}$, the no-load current
$I_c=\omega C_1 \ell\, U_{\mathrm{ph}}$ must equal the analytic value.
\item \textbf{Step / reflected wave.} For a fast voltage edge at the sending end,
the receiving-end overshoot and its rise must approach the lossless prediction set
by $\tau$ and the terminal reflection coefficient
$\Gamma=(Z_L-Z_0)/(Z_L+Z_0)$ as $N$ increases.
\end{enumerate}
Convergence of these three indicators with increasing $N$ confirms that the chosen
section count is sufficient.

%======================================================================
\section{Application notes}
%======================================================================
The same parameters quantify two phenomena common in converter-fed installations:

\paragraph{Reflected-wave overvoltage.} With $Z_0\approx20\ \Omega$ and a motor
input impedance of hundreds of ohms, $\Gamma\to1$ and the receiving-end voltage can
approach twice the DC-link-derived pulse. The critical cable length above which
full doubling occurs is $\ell_c=v\,t_r/2$, with $t_r$ the converter edge rise
time; for $t_r=100\ \mathrm{ns}$, $\ell_c\approx4\ \mathrm{m}$, so the
$120\ \mathrm{m}$ reference cable is firmly in the long-cable regime.

\paragraph{Common-mode current.} The $3C_g=108\ \mathrm{nF}$ shield path, excited
by the converter common-mode voltage, sets the magnitude and spectrum of the
common-mode current returning through the shield bonds. Single-point versus
both-point bonding is studied directly by changing the shield termination in the
model.

%======================================================================
\section*{References}
\addcontentsline{toc}{section}{References}
%======================================================================
\begin{thebibliography}{99}

\bibitem{paul} C.~R.~Paul, \emph{Analysis of Multiconductor Transmission Lines},
2nd ed. Hoboken, NJ: Wiley--IEEE Press, 2008.

\bibitem{marti} J.~R.~Mart\'i, ``Accurate modelling of frequency-dependent
transmission lines in electromagnetic transient simulations,'' \emph{IEEE Trans.
Power App. Syst.}, vol.~PAS-101, no.~1, pp.~147--157, Jan.~1982.

\bibitem{gustavsen} B.~Gustavsen and A.~Semlyen, ``Rational approximation of
frequency domain responses by vector fitting,'' \emph{IEEE Trans. Power Del.},
vol.~14, no.~3, pp.~1052--1061, Jul.~1999.

\bibitem{vonjouanne} A.~von Jouanne and P.~N.~Enjeti, ``Design considerations for
an inverter output filter to mitigate the effects of long motor leads in ASD
applications,'' \emph{IEEE Trans. Ind. Appl.}, vol.~33, no.~5, pp.~1138--1145,
Sep./Oct.~1997.

\bibitem{moreira} A.~F.~Moreira, T.~A.~Lipo, G.~Venkataramanan, and S.~Bernet,
``High-frequency modeling for cable and induction motor overvoltage studies in
long cable drives,'' \emph{IEEE Trans. Ind. Appl.}, vol.~38, no.~5,
pp.~1297--1306, Sep./Oct.~2002.

\bibitem{wang} L.~Wang, C.~N.~M.~Ho, F.~Canales, and J.~Jatskevich,
``High-frequency modeling of the long-cable-fed induction motor drive system using
TLM approach,'' \emph{IEEE Trans. Power Electron.}, vol.~25, no.~10,
pp.~2653--2664, Oct.~2010.

\bibitem{skibinski} G.~L.~Skibinski, R.~J.~Kerkman, and D.~Schlegel, ``EMI emissions
of modern PWM AC drives,'' \emph{IEEE Ind. Appl. Mag.}, vol.~5, no.~6,
pp.~47--80, Nov./Dec.~1999.

\bibitem{ieee575} \emph{IEEE Guide for Bonding Shields and Sheaths of
Single-Conductor Power Cables Rated 5~kV through 500~kV}, IEEE Std~575-2014, 2014.

\bibitem{iec60287} \emph{Electric Cables --- Calculation of the Current Rating ---
Part 1-1: Current Rating Equations (100\,\% Load Factor) and Calculation of
Losses}, IEC~60287-1-1, 2014.

\end{thebibliography}

%======================================================================
\appendix
\section{Parameter summary for Simscape}
%======================================================================
\begin{table}[h]
\centering
\caption{Direct entries for the Simscape three-phase $\pi$-line block
($\ell=0.12\ \mathrm{km}$).}
\begin{tabular}{@{}lcc@{}}
\toprule
Field & Positive sequence & Zero sequence \\
\midrule
$R$ ($\Omega/\mathrm{km}$) & $0.248$ & $0.248 + 3R_s' = 3.49$\,\footnotemark \\
$L$ ($\mathrm{mH/km}$) & $0.25$ & $0.27$ \\
$C$ ($\mu\mathrm{F/km}$) & $0.60$ & $0.30$ \\
Length (km) & \multicolumn{2}{c}{$0.12$} \\
Sections $N$ & \multicolumn{2}{c}{from Table~\ref{tab:N}} \\
\bottomrule
\end{tabular}
\end{table}
\footnotetext{The zero-sequence resistance includes the shield/earth return
contribution; the exact value depends on the bonding scheme and is finalized after
the Kron reduction \eqref{eq:kron} for the chosen installation.}

\end{document}
